%! Right Angles Revisited --- Setting Up Orthogonality — Unit 4, Lesson 4.0 — Guided Notes
\documentclass[10pt]{article}
\usepackage{linalg-article}
\usepackage{linalg-boxes}
\usetikzlibrary{calc}
\newcommand{\TallMath}[1]{$\displaystyle #1\rule[-1.4em]{0pt}{3.2em}$}
\newcommand{\vv}{\mathbf{v}}
\newcommand{\ww}{\mathbf{w}}
\newcommand{\avec}{\mathbf{a}}
\newcommand{\bb}{\mathbf{b}}
\newcommand{\xx}{\mathbf{x}}
\newcommand{\pp}{\mathbf{p}}
\newcommand{\zero}{\mathbf{0}}
\newcommand{\T}{\mathsf{T}}

\begin{document}

\pageheader{Unit 4, Lesson 4.0}{Guided Notes}
\namedateperiod

\vspace{0.12in}

\begin{objectivebox}
\textbf{By the end of this lesson, I will be able to\dots}
\begin{itemize}\setlength{\itemsep}{2pt}
  \item compute a \emph{dot product}, a \emph{length} $\|\vv\|=\sqrt{\vv\cdot\vv}$, and a \emph{unit vector};
  \item test whether two vectors are \emph{perpendicular} with $\vv\cdot\ww=0$, and solve for a missing entry;
  \item explain why the four subspaces are \emph{perpendicular pairs} and where Unit~4 is headed.
\end{itemize}
\end{objectivebox}

\vspace{0.08in}

\begin{vocabbox}
\small
Fill in each term as we recall or build it together.
\vspace{2pt}
\termblanklong{Dot product $\vv\cdot\ww$}
\termblanklong{Length $\|\vv\|$ and unit vector}
\termblanklong{Orthogonal (perpendicular)}
\termblanklong{Orthogonal complement (recall §3.6)}
\termblanklong{Projection (preview)}
\end{vocabbox}

\begin{hookbox}
A hiker stands \emph{off} a straight trail and wants the shortest walk back to it.
\begin{itemize}\setlength{\itemsep}{1pt}
  \item Of all the routes from the hiker back to the trail, which is \emph{shortest}? \writeline
  \item What is the \emph{angle} between that shortest route and the trail? \writeline
\end{itemize}
\end{hookbox}

\begin{notesbox}{1. Recall: Dot Product, Length, Unit Vector}
From Lesson~1.2, the \textbf{dot product} multiplies matching components and adds:
\[
  \vv\cdot\ww = v_1w_1 + v_2w_2 + \cdots + v_nw_n \quad(\text{one \blank{2.2cm}}).
\]
Its \textbf{sign} reports the angle: positive $=$ \blank{2.4cm}, zero $=$ a \blank{2.4cm} angle,
negative $=$ obtuse.

\vspace{3pt}
The \textbf{length} (or norm) is $\|\vv\|=\sqrt{\vv\cdot\vv}$. For
$\vv=\begin{bsmallmatrix} 3\\4 \end{bsmallmatrix}$:
$\|\vv\|=\sqrt{3^2+4^2}=\blank{1.2cm}$. To make a \textbf{unit vector} (same direction, length~1),
\blank{2.0cm} by the length: $\dfrac{\vv}{\|\vv\|}=\begin{bsmallmatrix} 3/5\\4/5 \end{bsmallmatrix}$.
\end{notesbox}

\begin{notesbox}{2. The One Big Idea: Perpendicular $\Leftrightarrow$ Dot $=0$}
Two nonzero vectors meet at a \textbf{right angle} exactly when
\[
  \vv\cdot\ww = \blank{1.2cm}.
\]
This single rule runs all of Unit~4: it turns the \emph{picture} of a right angle into an
\emph{equation}. \textbf{Check} the pair on the right:
$\begin{bsmallmatrix} 3\\4 \end{bsmallmatrix}\cdot\begin{bsmallmatrix} 4\\-3 \end{bsmallmatrix}
= 12 + (\blank{1.0cm}) = 0$, so they are \blank{2.6cm}.

\begin{center}
\begin{tikzpicture}[scale=0.5]
  \draw[step=1, linegray!45, thin] (-3.4,-3.4) grid (4.4,4.4);
  \draw[->, thick, charcoal] (-3.6,0)--(4.6,0);
  \draw[->, thick, charcoal] (0,-3.6)--(0,4.6);
  \draw[->, very thick, burgundy] (0,0)--(3,4) node[above right]{\scriptsize $\begin{bsmallmatrix} 3\\4 \end{bsmallmatrix}$};
  \draw[->, very thick, royalblue] (0,0)--(4,-3) node[below right]{\scriptsize $\begin{bsmallmatrix} 4\\-3 \end{bsmallmatrix}$};
  % right-angle mark at the origin between the two arrows
  \draw[charcoal, thin] (0.48,0.64)--(0.88,0.34)--(0.58,-0.06);
\end{tikzpicture}
\end{center}

\textbf{You can also solve for orthogonality.} What value of $x$ makes
$\begin{bsmallmatrix} 2\\3 \end{bsmallmatrix}$ perpendicular to
$\begin{bsmallmatrix} x\\-2 \end{bsmallmatrix}$? Set the dot product to zero:
$2x + 3(-2) = 0$, so $2x = 6$ and $x = \blank{1.2cm}$.
\end{notesbox}

\begin{notesbox}{3. The Four Subspaces Are Perpendicular Pairs (Recall §3.6)}
We ended Unit~3 with the big picture: the four fundamental subspaces come in two
\textbf{perpendicular pairs}. In the input space $\mathbb{R}^n$, the nullspace $N(A)$ is
perpendicular to the row space $C(A^{\T})$; in the output space, $N(A^{\T})\perp C(A)$.

\vspace{3pt}
\emph{Why} perpendicular? Read $A\xx=\zero$ \textbf{row by row}: each row of $A$ dotted with $\xx$
gives $0$. So every row is \blank{3.0cm} to $\xx$ --- that is just $\vv\cdot\ww=0$ again. Try it on
$A=\begin{bsmallmatrix} 1 & 2\\ 2 & 4 \end{bsmallmatrix}$: the nullspace is all multiples of
$\xx=\begin{bsmallmatrix} -2\\1 \end{bsmallmatrix}$. Check row~1:
$\begin{bsmallmatrix} 1\\2 \end{bsmallmatrix}\cdot\begin{bsmallmatrix} -2\\1 \end{bsmallmatrix}
= -2 + 2 = \blank{1.0cm}$. \checkmark\ Orthogonality is the engine of the whole picture.
\end{notesbox}

\begin{notesbox}{4. The Road Ahead: Drop a Perpendicular}
Back to the hiker. When a target $\bb$ is \emph{out of reach} of a line (or subspace), the
\textbf{closest} reachable point $\bb$ can get to is the \textbf{projection} $\pp$ --- the foot of
the perpendicular. The gap $\bb-\pp$ (the \emph{error}) is \blank{3.0cm} to the line.

\begin{center}
\begin{tikzpicture}[scale=1.0]
  \draw[->, very thick, burgundy] (-0.6,-0.4)--(4.6,3.07) node[right]{\scriptsize line (reachable)};
  \coordinate (O) at (0,0);
  \coordinate (b) at (1.4,3.0);
  \coordinate (p) at (2.53,1.69);   % foot of perpendicular from b to the line y=(2/3)x
  \fill[charcoal] (b) circle (1.6pt) node[above left]{\scriptsize target $\bb$};
  \fill[charcoal] (p) circle (1.6pt) node[below right]{\scriptsize projection $\pp$};
  \draw[->, very thick, royalblue] (p)--(b) node[midway, above left]{\scriptsize error $\bb-\pp$};
  % right-angle mark at p
  \draw[charcoal, thin] ($(p)+(-0.20,0.30)$)--($(p)+(-0.44,0.14)$)--($(p)+(-0.28,-0.10)$);
\end{tikzpicture}
\end{center}

\textbf{The road through Unit 4.} \textbf{4.1} orthogonality of the four subspaces (the §3.6 pairs,
in full)\; $\bullet$\; \textbf{4.2} \emph{projections} onto lines and subspaces (the closest
point)\; $\bullet$\; \textbf{4.3} \emph{least squares} --- best-fit lines when no exact answer
exists\; $\bullet$\; \textbf{4.4} orthogonal matrices $Q$ and Gram--Schmidt. Every step is the same
move: \emph{make something} \blank{2.6cm}.
\end{notesbox}

\begin{practicebox}
Use $\vv\cdot\ww=0$ for perpendicular. Show your dot products.
\begin{enumerate}[leftmargin=1.6em, itemsep=4pt]
  \item Are $\begin{bsmallmatrix} 5\\2 \end{bsmallmatrix}$ and $\begin{bsmallmatrix} 2\\-5 \end{bsmallmatrix}$ perpendicular? \writeline
  \item Find the length and unit vector of $\begin{bsmallmatrix} 6\\8 \end{bsmallmatrix}$. \writeline
  \item Find $x$ so that $\begin{bsmallmatrix} 4\\x \end{bsmallmatrix}$ is perpendicular to $\begin{bsmallmatrix} 3\\6 \end{bsmallmatrix}$. \writeline
\end{enumerate}
\end{practicebox}

\end{document}
